% Generated draft: requires mathematical and source review.
\documentclass[11pt]{article}
\usepackage[margin=25mm]{geometry}
\usepackage{fontspec}
\setmainfont{Latin Modern Roman}
\usepackage{amsmath,amssymb,mathtools,unicode-math}
\setmathfont{Latin Modern Math}
\usepackage{enumitem,needspace,xurl,hyperref}
\setlength{\parindent}{0pt}
\setlength{\parskip}{5pt}
\setlength{\emergencystretch}{3em}
\newcommand{\N}{\mathbb{N}}
\newcommand{\Z}{\mathbb{Z}}
\newcommand{\Q}{\mathbb{Q}}
\newcommand{\R}{\mathbb{R}}
\newcommand{\C}{\mathbb{C}}
\newcommand{\F}{\mathbb{F}}
\newcommand{\A}{\mathbb{A}}
\newcommand{\PP}{\mathbb{P}}
\newcommand{\E}{\mathbb{E}}
\newcommand{\eps}{\varepsilon}
\newcommand{\dd}{\,\mathrm{d}}
\begin{document}
\section*{UnsolvedMath ID: 800458}
\textbf{Source problem code:} catalogue entry
\subsection{Definitions and mathematical statement}
Let \((M_i^n,g_i,p_i)\) be pointed Riemannian manifolds of common dimension \(n\), and let \((X,d,p)\) be the limit appearing in the supplied statement. Write \(B_1^{g_i}(p_i)=\{q\in M_i:d_{g_i}(q,p_i)<1\}\). The stated hypotheses are
\[
\operatorname{Ric}_{g_i}\geq -(n-1)g_i,
\qquad
\operatorname{Vol}_{g_i}\bigl(B_1^{g_i}(p_i)\bigr)>v
\quad\text{for every }i,
\]
where \(v>0\) is independent of \(i\).

A tangent cone at \(x\in X\) is a pointed metric space \((Y,d_Y,y)\) obtained as a pointed Gromov--Hausdorff limit
\[
(X,r_j^{-1}d,x)\longrightarrow (Y,d_Y,y),
\qquad r_j>0,\quad r_j\downarrow0.
\]
Uniqueness means uniqueness up to an isometry preserving the distinguished point, independently of the sequence of scales and the convergent subsequence.

\textbf{Conjecture attributed to Naber (2020), Conjecture 2.} Under the stated hypotheses, the tangent cones of \(X\) are unique away from a set of codimension three.

\textbf{Qualifications needed for a precise formulation.} The supplied statement does not specify the convergence denoted by \((M_i^n,g_i,p_i)\to(X,d,p)\), whether the manifolds are complete or have boundary, or the meaning of codimension. Under the conventional interpretation of the convergence as pointed Gromov--Hausdorff convergence, and if ``a set of codimension three'' means an exceptional set of Hausdorff dimension at most \(n-3\), the intended conclusion can be written as follows: there exists \(E\subseteq X\) such that
\[
\dim_{\mathrm H}E\leq n-3
\]
and, for every \(x\in X\setminus E\), at least one tangent cone exists and any two tangent cones at \(x\) are pointed isometric. These interpretations require confirmation from the cited source. The supplied record gives no restriction on \(n\) or convention for the negative dimension bound when \(n<3\). Completeness and boundary hypotheses remain unresolved rather than being added to the conjecture.
\subsection{Short English statement}
Naber's supplied conjecture predicts that a limit of \(n\)-dimensional manifolds with the stated uniform Ricci lower bound and positive unit-ball volume lower bound has a unique tangent cone at every point outside an exceptional set of codimension three. The convergence, manifold hypotheses, and meaning of codimension require confirmation. The registry labels the problem open; this draft does not independently verify that status.
\subsection{Sources}
\begin{itemize}
\item Local UnsolvedMath registry, supplied record with permanent ID \(800458\), exact source problem code \texttt{catalogue entry}, title ``Spaces with Lower Ricci Curvature --- Conjecture 2.'' The record attributes the conjecture to Naber, with proposed year 2020, and labels its status \texttt{open}.
\item Supplied source reference: Naber, \emph{Conjectures and Open Questions on Spaces with Lower Ricci Curvature} (2020), Conjecture 2, \url{https://arxiv.org/abs/2010.10031}. The record reports access on 2026-10-16 and extraction method \texttt{arxiv\_latex}; these are supplied provenance statements, not independent verification or a claim of source access for this draft.
\item Supplied upstream revision: \texttt{d371ee917ac9647fbb06cd6fc529deec6f3dda82}.
\end{itemize}
\end{document}
